\documentclass[11pt,a4paper]{article}
\usepackage[utf8]{inputenc}
\usepackage[T1]{fontenc}
\usepackage[spanish]{babel}
\usepackage[margin=2.5cm]{geometry}
\usepackage{booktabs}
\usepackage{array}
\usepackage{tabularx}
\usepackage{listings}
\usepackage{xcolor}
\usepackage{enumitem}
\usepackage{parskip}
\usepackage{fancyhdr}
\usepackage{amsmath}
\usepackage{amssymb}

\definecolor{codegray}{RGB}{245,245,245}
\definecolor{codeblue}{RGB}{0,60,180}
\definecolor{codegreen}{RGB}{0,120,0}
\definecolor{codepurple}{RGB}{120,0,120}

\lstdefinestyle{sql}{
  language=SQL,
  backgroundcolor=\color{codegray},
  basicstyle=\ttfamily\small,
  keywordstyle=\color{codeblue}\bfseries,
  commentstyle=\color{codegreen},
  stringstyle=\color{codepurple},
  breaklines=true,
  frame=single,
  framerule=0.4pt,
  xleftmargin=8pt,
  xrightmargin=8pt,
  numbers=none,
  tabsize=2,
  showstringspaces=false,
  morekeywords={FETCH,FIRST,ROWS,ONLY,OFFSET,OVER,PARTITION,DENSE_RANK,ROW_NUMBER,NTILE,LAG,LEAD,FIRST_VALUE,LAST_VALUE,INTERSECT,EXCEPT}
}
\lstset{style=sql}

\pagestyle{fancy}
\fancyhf{}
\rhead{\textcolor{gray}{\small TD7: Ingeniería de Datos — UTDT}}
\lhead{\textcolor{gray}{\small Primer Parcial Modelo 1}}
\rfoot{\thepage}
\renewcommand{\headrulewidth}{0.4pt}

\setlist[enumerate,1]{label=\textbf{\alph*)}}
\setlist[enumerate,2]{label=\roman*.}

\begin{document}

\begin{center}
  {\LARGE\bfseries TD7 — Ingeniería de Datos}\\[6pt]
  {\large Universidad Torcuato Di Tella}\\[4pt]
  {\Large\bfseries Primer Parcial — Modelo 1}\\[4pt]
  {\normalsize Duración: 2 horas \quad|\quad Puntaje total: 100 puntos}
\end{center}
\noindent\rule{\linewidth}{0.8pt}
\vspace{2pt}

\noindent\textbf{Instrucciones:} Responda todos los ejercicios. Justifique cada respuesta. No está permitido el uso de material de consulta ni dispositivos electrónicos. Indique claramente el número de ejercicio y subejercicio.

\vspace{6pt}
\noindent\rule{\linewidth}{0.4pt}

% ─────────────────────────────────────────────────────────────
\section*{Ejercicio 1 — Modelo Entidad-Interrelación y Modelo Relacional \hfill\normalfont(30 puntos)}

Una plataforma de streaming de música necesita modelar su catálogo. Las reglas de negocio son:

\begin{itemize}
  \item Un \textbf{artista} tiene un identificador único, un nombre artístico y un año de debut.
  \item Un \textbf{álbum} tiene un código, un título y un año de lanzamiento. Todo álbum pertenece a \emph{exactamente un} artista. Un artista puede tener varios álbumes o ninguno.
  \item Cada \textbf{canción} tiene un identificador, un título y una duración en segundos. Una canción pertenece a \emph{exactamente un} álbum. Todo álbum tiene \emph{al menos una} canción.
  \item Los \textbf{usuarios} tienen un identificador, nombre, email y fecha de registro. Un usuario puede crear una o más \textbf{listas de reproducción}; cada lista tiene un nombre y pertenece a \emph{un único} usuario. Una canción puede estar en varias listas, y una lista puede contener varias canciones.
  \item Un usuario puede marcar canciones como \textbf{favoritas}. El sistema registra la fecha en que cada canción fue marcada favorita por cada usuario.
\end{itemize}

\begin{enumerate}
  \item (20 pts) Dibuje el DER completo. Indique entidades, atributos (marcando el atributo identificador de cada entidad), interrelaciones, cardinalidades y tipos de participación (total o parcial).
  \vspace{10cm}
  \item (10 pts) Traduzca el DER al modelo relacional. Para cada tabla indique el nombre de la relación, todos sus atributos, y marque con \underline{subrayado} la clave primaria y con cursiva las claves foráneas (indicando a qué tabla referencian).
\end{enumerate}

\newpage

% ─────────────────────────────────────────────────────────────
\section*{Ejercicio 2 — Álgebra Relacional \hfill\normalfont(25 puntos)}

Se tienen las siguientes relaciones del Mundial de Fútbol 2022:

\begin{center}
\begin{tabular}{ll}
\toprule
\textbf{Relación} & \textbf{Atributos} \\
\midrule
\texttt{Players} & \texttt{id, name, birth\_year, playing\_position, national\_team} \\
\texttt{Matches}  & \texttt{id, home, away, datetime, stage} \\
\texttt{Stages}   & \texttt{id, name} \\
\midrule
\end{tabular}
\end{center}

\noindent\small\textit{Notas: \texttt{home} y \texttt{away} son códigos de selección (ej. \texttt{'ARG'}). \texttt{stage} en \texttt{Matches} referencia a \texttt{Stages.id}.}

\normalsize
\vspace{4pt}
Escriba expresiones en Álgebra Relacional para cada consulta. Use la notación vista en clase ($\sigma$, $\pi$, $\cup$, $-$, $\cap$, $\times$, $\bowtie$, $\rho$).

\begin{enumerate}
  \item (4 pts) Obtener el \texttt{name} y \texttt{playing\_position} de todos los jugadores de Brasil (\texttt{national\_team = 'BRA'}).

  \vspace{2.2cm}

  \item (4 pts) Obtener los \texttt{id} de los partidos disputados en la etapa con \texttt{id = 3}.

  \vspace{2.2cm}

  \item (5 pts) Obtener los \texttt{name} de los jugadores nacidos antes de 1995 que se desempeñan como arqueros (\texttt{playing\_position = 'GK'}).

  \vspace{2.5cm}

  \item (6 pts) Para los partidos en que Argentina (\texttt{home = 'ARG'}) jugó de local, obtener el \texttt{name} de cada jugador argentino y el \texttt{name} de la etapa correspondiente a ese partido. \emph{(Sugerencia: relacione \texttt{Players}, \texttt{Matches} y \texttt{Stages}.)}

  \vspace{3cm}

  \item (6 pts) Obtener los \texttt{name} de los jugadores que pertenecen a selecciones que \textbf{no} disputaron ningún partido (ni como local ni como visitante). \emph{(Sugerencia: use diferencia de conjuntos y renombramiento $\rho$.)}

  \vspace{3.5cm}
\end{enumerate}

\newpage

% ─────────────────────────────────────────────────────────────
\section*{Ejercicio 3 — SQL \hfill\normalfont(45 puntos)}

Use la base de datos \textbf{Chinook}. El esquema relevante es:

\begin{center}\small
\begin{tabular}{ll}
\toprule
\textbf{Tabla} & \textbf{Atributos principales} \\
\midrule
\texttt{Artist}       & \texttt{ArtistId (PK), Name} \\
\texttt{Album}        & \texttt{AlbumId (PK), Title, ArtistId (FK)} \\
\texttt{Track}        & \texttt{TrackId (PK), Name, AlbumId (FK), GenreId (FK), Composer, Milliseconds, UnitPrice} \\
\texttt{Genre}        & \texttt{GenreId (PK), Name} \\
\texttt{Invoice}      & \texttt{InvoiceId (PK), CustomerId (FK), InvoiceDate, BillingCountry, Total} \\
\texttt{InvoiceLine}  & \texttt{InvoiceLineId (PK), InvoiceId (FK), TrackId (FK), UnitPrice, Quantity} \\
\texttt{Customer}     & \texttt{CustomerId (PK), FirstName, LastName, SupportRepId (FK a Employee)} \\
\texttt{Employee}     & \texttt{EmployeeId (PK), FirstName, LastName, ReportsTo (FK a Employee)} \\
\texttt{Playlist}     & \texttt{PlaylistId (PK), Name} \\
\texttt{PlaylistTrack}& \texttt{PlaylistId (FK), TrackId (FK)} \\
\bottomrule
\end{tabular}
\end{center}

\begin{enumerate}

  \item (7 pts) Listar el \texttt{Name} de cada género musical junto con la cantidad de canciones que contiene. Incluir los géneros sin canciones (mostrando 0). Ordenar de mayor a menor cantidad.

  \vspace{4cm}

  \item (8 pts) Listar el nombre de los artistas que tienen \textbf{más de 10 álbumes}. Mostrar el nombre del artista y la cantidad de álbumes.

  \vspace{4cm}

  \item (8 pts) Listar el \texttt{FirstName} y \texttt{LastName} de los clientes que \textbf{nunca} realizaron ninguna compra (sin facturas). Utilizar un \texttt{LEFT JOIN}.

  \vspace{4cm}

  \item (8 pts) Para cada país (\texttt{BillingCountry}), calcular el monto total de ventas (suma de \texttt{Total}). Mostrar únicamente los países cuyo total supere \$100, ordenados de mayor a menor.

  \vspace{4cm}

  \item (14 pts) Usando una \textbf{CTE}, encontrar la playlist con mayor cantidad de canciones. Mostrar el nombre de la playlist y la cantidad de canciones. \emph{Si hay empate, mostrar todas las que comparten el máximo.}

  \vspace{5cm}

\end{enumerate}

\end{document}
